% ===========================================================================================
% \chapter{PREAMBLE}

\documentclass[12pt, a4paper]{book}

% ── Encoding & fonts (XeLaTeX) ───────────────────────────────────
\usepackage{fontspec}
\setmainfont{Times New Roman}
\setsansfont{Arial}
\setmonofont{Courier New}

% ── Language ─────────────────────────────────────────────────────
\usepackage{polyglossia}
\setdefaultlanguage{english}

% ── Mathematics ──────────────────────────────────────────────────
\usepackage{amsmath}
\usepackage{amssymb}
\usepackage{mathtools}

% ── Page layout ──────────────────────────────────────────────────
\usepackage[
  a4paper,
  top    = 2.5cm,
  bottom = 2.5cm,
  left   = 3.0cm,
  right  = 2.5cm
]{geometry}

% ── Tables ───────────────────────────────────────────────────────
\usepackage{booktabs}
\usepackage{array}
\usepackage{longtable}

% ── Figures ──────────────────────────────────────────────────────
\usepackage{graphicx}
\usepackage{float}

% ── Hyperlinks ───────────────────────────────────────────────────
\usepackage[
  colorlinks = true,
  linkcolor  = blue,
  citecolor  = blue,
  urlcolor   = blue
]{hyperref}

% ── Index (MakeIndex) ────────────────────────────────────────────
\usepackage{makeidx}
\makeindex

% ── Bibliography (BIBTeX) ────────────────────────────────────────
\usepackage{natbib}

% ── Units ────────────────────────────────────────────────────────
\usepackage{siunitx}
\sisetup{output-decimal-marker={.}}

% ── Source code listing ──────────────────────────────────────────
\usepackage{listings}
\lstset{
  basicstyle   = \ttfamily\small,
  breaklines   = true,
  frame        = single,
  numbers      = left,
  numberstyle  = \tiny,
  keywordstyle = \bfseries,
  language     = Python
}

% ── Custom colours ───────────────────────────────────────────────
\usepackage{xcolor}
\definecolor{pyblue}{HTML}{306998}

% ===========================================================================================
% \subsection{Image}
% ===========================================================================================
\usepackage{graphicx}
\graphicspath{
	{../images/}
	{}
	{}
}

% ─────────────────────────────────────────────────────────────────
%  AUTO-GENERATED VALUES from lines_plan.ipynb
%  File ini ditulis oleh notebook; jangan edit manual.
% ─────────────────────────────────────────────────────────────────
% ============================================================
% AUTO-GENERATED oleh lines_plan.py -- 2026-05-29 18:00:55
% Sumber Excel : Lines_Plan_-_Oil_Tanker_20704_DWT_-_Manual.xlsx
% Sheet        : Lines Plan, sel B434:G458
% JANGAN EDIT MANUAL -- jalankan lines_plan.py untuk regenerasi.
% ============================================================

% -- Nilai global --------------------------------------------
\newcommand{\LPSumPct}{1.3576}
\newcommand{\LPJumlahStation}{23}

% -- Nilai per station ---------------------------------------
\newcommand{\LPStaMinTwoAone}{0.0000}
\newcommand{\LPStaMinTwoAtwo}{0.0000}
\newcommand{\LPStaMinTwoSel}{0.000000}
\newcommand{\LPStaMinTwoPct}{0.0000}

\newcommand{\LPStaMinOneAone}{0.0000}
\newcommand{\LPStaMinOneAtwo}{0.0000}
\newcommand{\LPStaMinOneSel}{0.000000}
\newcommand{\LPStaMinOnePct}{0.0000}

\newcommand{\LPStaZeroAone}{0.0000}
\newcommand{\LPStaZeroAtwo}{0.0000}
\newcommand{\LPStaZeroSel}{0.000000}
\newcommand{\LPStaZeroPct}{0.0000}

\newcommand{\LPStaOneAone}{4.2867}
\newcommand{\LPStaOneAtwo}{4.2935}
\newcommand{\LPStaOneSel}{0.006800}
\newcommand{\LPStaOnePct}{0.1586}

\newcommand{\LPStaTwoAone}{2.6553}
\newcommand{\LPStaTwoAtwo}{2.6531}
\newcommand{\LPStaTwoSel}{0.002200}
\newcommand{\LPStaTwoPct}{0.0829}

\newcommand{\LPStaThreeAone}{2.3475}
\newcommand{\LPStaThreeAtwo}{2.3470}
\newcommand{\LPStaThreeSel}{0.000500}
\newcommand{\LPStaThreePct}{0.0213}

\newcommand{\LPStaFourAone}{0.9877}
\newcommand{\LPStaFourAtwo}{0.9880}
\newcommand{\LPStaFourSel}{0.000300}
\newcommand{\LPStaFourPct}{0.0304}

\newcommand{\LPStaFiveAone}{0.3218}
\newcommand{\LPStaFiveAtwo}{0.3217}
\newcommand{\LPStaFiveSel}{0.000100}
\newcommand{\LPStaFivePct}{0.0311}

\newcommand{\LPStaSixAone}{0.4044}
\newcommand{\LPStaSixAtwo}{0.4046}
\newcommand{\LPStaSixSel}{0.000200}
\newcommand{\LPStaSixPct}{0.0495}

\newcommand{\LPStaSevenAone}{0.1988}
\newcommand{\LPStaSevenAtwo}{0.1987}
\newcommand{\LPStaSevenSel}{0.000100}
\newcommand{\LPStaSevenPct}{0.0503}

\newcommand{\LPStaEightAone}{0.1988}
\newcommand{\LPStaEightAtwo}{0.1987}
\newcommand{\LPStaEightSel}{0.000100}
\newcommand{\LPStaEightPct}{0.0503}

\newcommand{\LPStaNineAone}{0.1988}
\newcommand{\LPStaNineAtwo}{0.1987}
\newcommand{\LPStaNineSel}{0.000100}
\newcommand{\LPStaNinePct}{0.0503}

\newcommand{\LPStaTenAone}{0.1988}
\newcommand{\LPStaTenAtwo}{0.1987}
\newcommand{\LPStaTenSel}{0.000100}
\newcommand{\LPStaTenPct}{0.0503}

\newcommand{\LPStaElevenAone}{0.1988}
\newcommand{\LPStaElevenAtwo}{0.1987}
\newcommand{\LPStaElevenSel}{0.000100}
\newcommand{\LPStaElevenPct}{0.0503}

\newcommand{\LPStaTwelveAone}{0.1988}
\newcommand{\LPStaTwelveAtwo}{0.1987}
\newcommand{\LPStaTwelveSel}{0.000100}
\newcommand{\LPStaTwelvePct}{0.0503}

\newcommand{\LPStaThirteenAone}{0.1988}
\newcommand{\LPStaThirteenAtwo}{0.1987}
\newcommand{\LPStaThirteenSel}{0.000100}
\newcommand{\LPStaThirteenPct}{0.0503}

\newcommand{\LPStaFourteenAone}{0.1988}
\newcommand{\LPStaFourteenAtwo}{0.1987}
\newcommand{\LPStaFourteenSel}{0.000100}
\newcommand{\LPStaFourteenPct}{0.0503}

\newcommand{\LPStaFifteenAone}{1.4231}
\newcommand{\LPStaFifteenAtwo}{1.4223}
\newcommand{\LPStaFifteenSel}{0.000800}
\newcommand{\LPStaFifteenPct}{0.0562}

\newcommand{\LPStaSixteenAone}{2.3770}
\newcommand{\LPStaSixteenAtwo}{2.3773}
\newcommand{\LPStaSixteenSel}{0.000300}
\newcommand{\LPStaSixteenPct}{0.0126}

\newcommand{\LPStaSeventeenAone}{2.5714}
\newcommand{\LPStaSeventeenAtwo}{2.5701}
\newcommand{\LPStaSeventeenSel}{0.001300}
\newcommand{\LPStaSeventeenPct}{0.0506}

\newcommand{\LPStaEighteenAone}{3.6467}
\newcommand{\LPStaEighteenAtwo}{3.6471}
\newcommand{\LPStaEighteenSel}{0.000400}
\newcommand{\LPStaEighteenPct}{0.0110}

\newcommand{\LPStaNineteenAone}{2.0589}
\newcommand{\LPStaNineteenAtwo}{2.0596}
\newcommand{\LPStaNineteenSel}{0.000700}
\newcommand{\LPStaNineteenPct}{0.0340}

\newcommand{\LPStaTwentyAone}{2.3251}
\newcommand{\LPStaTwentyAtwo}{2.3348}
\newcommand{\LPStaTwentySel}{0.009700}
\newcommand{\LPStaTwentyPct}{0.4172}


% ── Document metadata ────────────────────────────────────────────
\title{%
  \textbf{Verifikasi Lines Plan --- Oil Tanker 20704 DWT}\\[0.4em]
  \large Laporan Perhitungan Otomatis dari Jupyter Notebook
}
\author{Daffa \\ \small Politeknik Perkapalan Negeri Surabaya}
\date{\today}

% ─────────────────────────────────────────────────────────────────
\begin{document}
% ─────────────────────────────────────────────────────────────────

\maketitle
\tableofcontents
\newpage

% ─────────────────────────────────────────────────────────────────
\section{Pendahuluan}
\label{sec:intro}
% ─────────────────────────────────────────────────────────────────

\index{Lines Plan}
\index{Oil Tanker}
Laporan ini menyajikan hasil verifikasi \emph{luas penampang} (\textit{cross-sectional area})
tiap \emph{station} dari rencana garis (\textit{Lines Plan}) kapal tanker minyak
berkapasitas \textbf{20.704 DWT}.

Nilai-nilai numerik dihitung dalam \texttt{lines\_plan.ipynb} dan diekspor secara otomatis
ke \texttt{calc\_values.tex} melalui mekanisme \verb|\newcommand|, sehingga
setiap angka dalam laporan ini \emph{selalu sinkron} dengan notebook --- tanpa
risiko salah salin (\textit{copy-paste error}).

% ─────────────────────────────────────────────────────────────────
\section{Metodologi Verifikasi}
\label{sec:method}
% ─────────────────────────────────────────────────────────────────

\index{verifikasi luas}
Verifikasi dilakukan dengan membandingkan dua hasil integrasi luas penampang:

\begin{itemize}
  \item $A_1$ --- luas hasil metode integrasi pertama (\si{m^2})
  \item $A_2$ --- luas hasil metode integrasi kedua (\si{m^2})
\end{itemize}

\noindent
Formula yang digunakan (mereplikasi formula Excel pada sel F dan G):

\begin{equation}
  \Delta = |A_1 - A_2|
  \label{eq:selisih}
\end{equation}

\begin{equation}
  \varepsilon\,[\%] = \frac{\Delta}{A_1} \times 100
  \label{eq:pct}
\end{equation}

\noindent
Total penyimpangan dihitung sebagai:

\begin{equation}
  \Sigma\,\varepsilon = \sum_{\text{semua station}} \varepsilon_i
  \label{eq:sum}
\end{equation}

% ─────────────────────────────────────────────────────────────────
\section{Hasil Perhitungan}
\label{sec:results}
% ─────────────────────────────────────────────────────────────────

\index{hasil perhitungan}
Tabel~\ref{tab:linesplan} menyajikan hasil verifikasi untuk seluruh
\LPJumlahStation\ station (station $-2$ s.d.\ $20$), yang mereplikasi
isi sel B434:G458 pada sheet \textit{Lines Plan}.

\begin{longtable}{crrrr}
  \caption{Verifikasi Luas Penampang per Station}
  \label{tab:linesplan} \\
  \toprule
  \textbf{Station} & \textbf{Luas $A_1$ (m$^2$)} & \textbf{Luas $A_2$ (m$^2$)} &
  \textbf{Selisih (m$^2$)} & \textbf{\% Selisih} \\
  \midrule
  \endfirsthead
  \multicolumn{5}{c}{\tablename\ \thetable{} (lanjutan)} \\
  \toprule
  \textbf{Station} & \textbf{Luas $A_1$ (m$^2$)} & \textbf{Luas $A_2$ (m$^2$)} &
  \textbf{Selisih (m$^2$)} & \textbf{\% Selisih} \\
  \midrule
  \endhead
  \bottomrule
  \endfoot
  $-2$ & \LPStaMinTwoAone  & \LPStaMinTwoAtwo  & \LPStaMinTwoSel  & --- \\
  $-1$ & \LPStaMinOneAone  & \LPStaMinOneAtwo  & \LPStaMinOneSel  & --- \\
  $0$  & \LPStaZeroAone   & \LPStaZeroAtwo   & \LPStaZeroSel   & --- \\
  $1$  & \LPStaOneAone   & \LPStaOneAtwo   & \LPStaOneSel   & \LPStaOnePct  \\
  $2$  & \LPStaTwoAone   & \LPStaTwoAtwo   & \LPStaTwoSel   & \LPStaTwoPct  \\
  $3$  & \LPStaThreeAone   & \LPStaThreeAtwo   & \LPStaThreeSel   & \LPStaThreePct  \\
  $4$  & \LPStaFourAone   & \LPStaFourAtwo   & \LPStaFourSel   & \LPStaFourPct  \\
  $5$  & \LPStaFiveAone   & \LPStaFiveAtwo   & \LPStaFiveSel   & \LPStaFivePct  \\
  $6$  & \LPStaSixAone   & \LPStaSixAtwo   & \LPStaSixSel   & \LPStaSixPct  \\
  $7$  & \LPStaSevenAone   & \LPStaSevenAtwo   & \LPStaSevenSel   & \LPStaSevenPct  \\
  $8$  & \LPStaEightAone   & \LPStaEightAtwo   & \LPStaEightSel   & \LPStaEightPct  \\
  $9$  & \LPStaNineAone   & \LPStaNineAtwo   & \LPStaNineSel   & \LPStaNinePct  \\
  $10$ & \LPStaTenAone  & \LPStaTenAtwo  & \LPStaTenSel  & \LPStaTenPct \\
  $11$ & \LPStaElevenAone  & \LPStaElevenAtwo  & \LPStaElevenSel  & \LPStaElevenPct \\
  $12$ & \LPStaTwelveAone  & \LPStaTwelveAtwo  & \LPStaTwelveSel  & \LPStaTwelvePct \\
  $13$ & \LPStaThirteenAone  & \LPStaThirteenAtwo  & \LPStaThirteenSel  & \LPStaThirteenPct \\
  $14$ & \LPStaFourteenAone  & \LPStaFourteenAtwo  & \LPStaFourteenSel  & \LPStaFourteenPct \\
  $15$ & \LPStaFifteenAone  & \LPStaFifteenAtwo  & \LPStaFifteenSel  & \LPStaFifteenPct \\
  $16$ & \LPStaSixteenAone  & \LPStaSixteenAtwo  & \LPStaSixteenSel  & \LPStaSixteenPct \\
  $17$ & \LPStaSeventeenAone  & \LPStaSeventeenAtwo  & \LPStaSeventeenSel  & \LPStaSeventeenPct \\
  $18$ & \LPStaEighteenAone  & \LPStaEighteenAtwo  & \LPStaEighteenSel  & \LPStaEighteenPct \\
  $19$ & \LPStaNineteenAone  & \LPStaNineteenAtwo  & \LPStaNineteenSel  & \LPStaNineteenPct \\
  $20$ & \LPStaTwentyAone  & \LPStaTwentyAtwo  & \LPStaTwentySel  & \LPStaTwentyPct \\
  \midrule
  \multicolumn{4}{r}{\textbf{SUM \% Selisih}} & \textbf{\LPSumPct} \\
  \bottomrule
\end{longtable}

\noindent
Total penyimpangan persentase seluruh station (replikasi sel G458):

\begin{equation}
  \Sigma\,\varepsilon = \LPSumPct\,\%
  \label{eq:sumresult}
\end{equation}

% ─────────────────────────────────────────────────────────────────
\section{Penjelasan Kode Python}
\label{sec:code}
% ─────────────────────────────────────────────────────────────────

\index{Jupyter Notebook}
\index{openpyxl}
Sel utama dalam \texttt{lines\_plan.ipynb} yang mereplikasi formula Excel:

\begin{lstlisting}[language=Python,
  caption={Replikasi formula Excel --- sel F dan G (Selisih dan \% Selisih)}]
for station, (A1, A2) in station_data.items():
    selisih = abs(A1 - A2)           # =ABS(D-E)
    pct = (selisih / A1 * 100) if A1 > 0 else None  # =(F/D)*100

# SUM % Selisih  -- replikasi =SUM(G435:G457)
sum_pct = sum(r["pct"] for r in results if r["pct"] is not None)
\end{lstlisting}

\noindent
Nilai kemudian diekspor ke \texttt{calc\_values.tex}:

\begin{lstlisting}[language=Python,
  caption={Ekspor nilai ke LaTeX via \textbackslash newcommand}]
lines += [
    rf"\newcommand{{\LP{mac}Aone}}{{{fmt(r['luas_A1'], 4)}}}",
    rf"\newcommand{{\LP{mac}Atwo}}{{{fmt(r['luas_A2'], 4)}}}",
    rf"\newcommand{{\LP{mac}Sel}}{{{fmt(r['selisih'], 6)}}}",
    rf"\newcommand{{\LP{mac}Pct}}{{{fmt(pct_val, 4)}}}",
]
\end{lstlisting}

% ─────────────────────────────────────────────────────────────────
\section{Kesimpulan}
\label{sec:conclusion}
% ─────────────────────────────────────────────────────────────────

Hasil verifikasi menunjukkan total penyimpangan luas penampang antar dua metode
integrasi sebesar $\Sigma\,\varepsilon = \LPSumPct\,\%$ untuk seluruh
\LPJumlahStation\ station (B434:G458).
Nilai ini konsisten dengan hasil pada file Excel asli (sel G458).

% ─────────────────────────────────────────────────────────────────
%  Bibliography (BIBTeX)
% ─────────────────────────────────────────────────────────────────
\nocite{*}
\bibliographystyle{plainnat}
\bibliography{references}

% ─────────────────────────────────────────────────────────────────
%  Index (MakeIndex)
% ─────────────────────────────────────────────────────────────────
\printindex

\end{document}